\documentclass[a4paper,12pt]{article}
\usepackage{graphicx}
\usepackage[utf8]{inputenc}  
\usepackage[T1]{fontenc}
\usepackage{gensymb}

\usepackage{tgheros,textcomp}%Police Arial
\renewcommand{\familydefault}{\sfdefault}

\usepackage[top = 1cm, bottom = 2cm, left = 1cm, right = 1cm]{geometry}
\usepackage{titlesec}
\setcounter{secnumdepth}{4}
\usepackage{xcolor}
\usepackage{amssymb} %double flèche de réaction
\usepackage{xtab}%tableau sur plusieurs pages
\usepackage{esvect}%grands vecteurs

\usepackage[version=4]{mhchem} %equations chimiques

\usepackage{array}
\usepackage{chemfig}
\usepackage{multirow}
\usepackage{sectsty}
\usepackage{caption}
\usepackage{amsmath}
\usepackage{amssymb}
\usepackage[cal=boondoxo]{mathalfa} 
\usepackage{enumitem}
\usepackage{lastpage}
\usepackage{tcolorbox}
\usepackage{siunitx}
\usepackage{tikz}
\usetikzlibrary{calc}
\usepackage{multicol}
\usepackage{eqnarray}

\sisetup{
	detect-all,
	output-decimal-marker={,},
	group-minimum-digits = 3,
	group-separator={~},
	number-unit-separator={~},
	inter-unit-product={~}
} %setup pour SIunitx

\usepackage{listings}%pour insérer le code python
\definecolor{codegreen}{rgb}{0,0.6,0}
\definecolor{codegray}{rgb}{0.5,0.5,0.5}
\definecolor{codepurple}{rgb}{0.58,0,0.82}
\definecolor{backcolour}{rgb}{0.95,0.95,0.92}

\graphicspath{ {/home/rweye/Documents/Enseignement/2022/2GT/C7/Ressources/} }

\sectionfont{\fontsize{14}{8}\selectfont}
\renewcommand{\thesection}{\arabic{section}.}
\renewcommand{\thesubsection}{\hspace{1cm}\alph{subsection}.}

\title{\vspace{-1cm}\large 
	\begin{tabular}{|m{7cm}|>{\centering}m{6cm}|>{\centering\arraybackslash}m{4cm}|}
		\hline
		 Nom :& \textbf{\textsc{Évaluation : Le son}} &  Chapitre 7\\
		 Prénom : &&\\
		\hline
\end{tabular}}
\date{}

\newpagestyle{mystyle}{%
	\setfoot{}{\thepage\,$|$\,\pageref{LastPage}}{}
}
\renewpagestyle{plain}{%
	\setfoot{}{\thepage\,$|$\,\pageref{LastPage}}{}
}

\pagestyle{mystyle}

\newenvironment{itemize*}%
{\begin{itemize}[label=$-$]%
		\setlength{\itemsep}{0pt}%
		\setlength{\parskip}{0pt}}%
	{\end{itemize}}

\newenvironment{itemize**}%
{\begin{itemize}[label=]%
		\setlength{\itemsep}{0pt}%
		\setlength{\parskip}{0pt}}%
	{\end{itemize}}

\begin{document}
	\maketitle
	\vspace{-3cm}
\section*{Exercice 1 : Sons}
\begin{tabular}{cc}
	\includegraphics[scale=.5]{2GT_C7_Eval_Fig1.png}&\includegraphics[scale=.5]{2GT_C7_Eval_Fig3.png}\\
	\includegraphics[scale=.5]{2GT_C7_Eval_Fig4.png}&\includegraphics[scale=.5]{2GT_C7_Eval_Fig2.png}\\
\end{tabular}
\begin{enumerate}
	\item Dire si les signaux sonores présentés ci dessus sont périodiques. Justifier chaque réponse.
	\item Pour un signal périodique, définir un motif. Repasser en couleur un motif pour chaque signal périodique trouvé dans la réponse précédente.
	\item Déterminer les fréquences des signaux précédent si possible (des calculs détaillés sont attendus à chaque fois).
\end{enumerate}

\section*{Exercice 2 : Écholocalisation}
L'écholocalisation est un système de perception de l'environnement développé par plusieurs espèces (chauve-souris, cétacés) afin de percevoir l'environnement proche dans utiliser la vision.\\
Il s'agit pour l'animal de produire un son de haute fréquence qui va se propager jusqu'à un obstacle puis revenir à l'oreille de l'animal. La perception de proximité est faite par la mesure du temps mis par l'onde sonore pour revenir à son émetteur.
\begin{enumerate}
	\item Proposer un schéma permettant de représenter le système d'écholocalisation. Faire figurer un émetteur ($E$), un objet à localiser ($A$), un récepteur ($R$) et la distance entre le système émetteur/récepteur et l'objet ($d$).
	\item Dans le cas d'une chauve-souris, exprimer le temps ($\Delta t_{CS}$)mis par l'onde sonore pour revenir vers elle après avoir rencontré un insecte situé à 3 m.
	\item Sachant que, le son a une célérité $c_{air}$ = \SI{340}{\meter\per\second}, calculer la valeur de $\Delta t_{CS}$.
	\item Dans le cas où l'animal est un dauphin, exprimer le temps ($\Delta t_{D}$)mis par l'onde sonore pour revenir vers elle après avoir rencontré un saumon situé à 3 m. Puis calculer $\Delta t_{D}$, sachant que $c_{eau}$ = \SI{1500}{\meter\per\second}
	\item Lequel des deux animaux a le potentiel d'être le plus rapide pour chasser ?
	\item \{Question supplémentaire non notée :\} imaginer d'autre paramètres pouvant influencer la vitesse en chasse des deux animaux
\end{enumerate}

\section*{Exercice 3 : Let's get loud !}

	Lors d'un concert, un groupe d'amis cherchent une place où ils pourront apprécier le spectacle dans les meilleures conditions. Trois avis se font concurrence :\\
	\noindent\begin{minipage}{.6\linewidth}
\begin{itemize}
	\item Avis 1 : "Devant ! On sera proche de la scène et on pourra voir les artistes de près et même avec un peu de chance récupérer des cadeaux !"
	\item Avis 2 : "Il vaudrait mieux être au milieu ... On pourra toujours voir les artistes, le son ne sera pas trop fort et on sera en plein dans l'ambiance pour danser !"
	\item Avis 3 : "Au fond, c'est plus proche de la sortie et des toilettes et on n'est pas chahuté par tout le monde."
\end{itemize}
\textbf{Document 2 : Données pour des enceintes de 200 W}\\
\begin{minipage}{\linewidth}
	\centering
	\includegraphics[scale=.5]{2GT_C7_Eval_Fig5.png}
\end{minipage}
\end{minipage}
\begin{minipage}{.4\linewidth}
	\centering
	\textbf{Document 1 :}\\
	\includegraphics[scale=.4]{2GT_C7_Eval_Fig6.png}
\end{minipage}

\textbf{Question : } En se basant sur l'énoncé et sur vos connaissances, justifier quel endroit du public vous jugez le plus sûr d'un point de vue du volume sonore.
\end{document}